% ===========================================================================
% LaTeX source generated by scripts/build_latex_papers.py (pandoc)
%
% JOURNAL TEMPLATE SWAP:
%   The default documentclass is "article". For final submission, replace the
%   \documentclass[...] line with the target journal's class. See COMPILE.md
%   for a per-journal template swap table and download instructions.
%
% COMPILE:
%   tectonic pN.tex              (single-binary, recommended)
%   OR: pdflatex pN.tex && bibtex pN && pdflatex pN.tex && pdflatex pN.tex
%
% REFERENCES:
%   references.bib is auto-extracted from the markdown References section.
%   Verify all metadata (author, title, journal, volume, pages, DOI) in
%   Zotero/JabRef before submission — stub entries contain the raw citation
%   text so nothing is lost.
%
% FIGURES:
%   Add real PNG/PDF figures under figures/ and update \includegraphics paths.
% ===========================================================================

\documentclass[
]{article}
\usepackage{xcolor}
\usepackage{amsmath,amssymb}
\setcounter{secnumdepth}{5}
\usepackage{iftex}
\ifPDFTeX
  \usepackage[T1]{fontenc}
  \usepackage[utf8]{inputenc}
  \usepackage{textcomp} % provide euro and other symbols
\else % if luatex or xetex
  \usepackage{unicode-math} % this also loads fontspec
  \defaultfontfeatures{Scale=MatchLowercase}
  \defaultfontfeatures[\rmfamily]{Ligatures=TeX,Scale=1}
\fi
\usepackage{lmodern}
\ifPDFTeX\else
  % xetex/luatex font selection
\fi
% Use upquote if available, for straight quotes in verbatim environments
\IfFileExists{upquote.sty}{\usepackage{upquote}}{}
\IfFileExists{microtype.sty}{% use microtype if available
  \usepackage[]{microtype}
  \UseMicrotypeSet[protrusion]{basicmath} % disable protrusion for tt fonts
}{}
\makeatletter
\@ifundefined{KOMAClassName}{% if non-KOMA class
  \IfFileExists{parskip.sty}{%
    \usepackage{parskip}
  }{% else
    \setlength{\parindent}{0pt}
    \setlength{\parskip}{6pt plus 2pt minus 1pt}}
}{% if KOMA class
  \KOMAoptions{parskip=half}}
\makeatother
\usepackage{longtable,booktabs,array}
\usepackage{caption}
\captionsetup[table]{skip=6pt}
\newcounter{none} % for unnumbered tables
\usepackage{calc} % for calculating minipage widths
% Correct order of tables after \paragraph or \subparagraph
\usepackage{etoolbox}
\makeatletter
\patchcmd\longtable{\par}{\if@noskipsec\mbox{}\fi\par}{}{}
\makeatother
\setlength{\emergencystretch}{3em} % prevent overfull lines
\providecommand{\tightlist}{%
  \setlength{\itemsep}{0pt}\setlength{\parskip}{0pt}}
\usepackage[]{natbib}
% Pandoc wrapper fallback (in case pandoc preamble is stripped)
\providecommand{\pandocbounded}[1]{#1}
\bibliographystyle{plainnat}
\usepackage{bookmark}
\IfFileExists{xurl.sty}{\usepackage{xurl}}{} % add URL line breaks if available
\urlstyle{same}
% fallback for those not using the hyperref driver hyperxmp:
\makeatletter
\@ifundefined{xmpquote}{\newcommand{\xmpquote}[1]{#1}}{}
\makeatother
\hypersetup{
  pdftitle={Preregistered Rejection as a First-Class Negative Result: A Methodological Case Study on the 116-Unit Drosophila Male CNS},
  pdfauthor={Li Xing},
  hidelinks,
  pdfcreator={LaTeX via pandoc}}

\title{Preregistered Rejection as a First-Class Negative Result: A
Methodological Case Study on the 116-Unit Drosophila Male CNS}
\author{Li Xing}
\date{2026-10-21}

\begin{document}
\maketitle
\begin{abstract}
We describe a methodological framework for preregistering, running, and
publishing computational experiments whose central contribution is
\emph{negative results} rather than positive effects. The framework
treats the preregistration document itself as a first-class research
object --- with its own version, its own failure modes, and its own
validation gate --- and makes three disciplines explicit: 1.
\textbf{Preregistration is time-ordered, machine-readable, and
one-shot.} Predictions, features, thresholds, and seeds are locked
\emph{before} execution and verified byte-for-byte at analysis time. 2.
\textbf{A preregistered \emph{rejection} is a valid endpoint.} If a
proposed analysis fails a pre-written acceptance rule, the rejection is
recorded as a result alongside the experiments that did run. 3.
\textbf{The measurement base has an upper bound.} When the available
metric family does not contain an independent second reading of an
intervention, the experiment is closed as \emph{base-limited}, not as
\emph{underpowered} or as \emph{not enough effort}. We illustrate the
framework on five preregistered protocols (PR-1, PR-2, PR-2b, PR-3,
PR-4). Four of the five are negative results; one is a preregistered
rejection of a proposed experiment that never ran. None of the five was
modified after the fact to flip a FAIL to a PASS. The full protocol, the
raw artifacts, and the three read-only metric scripts that reproduce the
paper's traceability figures are in \texttt{metrics\_fork/} at the root
of the upstream repository; a reader who pulls the upstream artifact and
runs the three commands gets the paper's numbers, or finds where we got
them wrong. \textbf{Keywords.} preregistration, negative results,
preregistered rejection, scientific self-verification, reproducibility,
base-limit honesty, computational neuroscience. ---
\end{abstract}

{
\setcounter{tocdepth}{3}
\tableofcontents
}
\section{Reproducible Preregistration as a First-Class Research Object:
A Methodology for Computing Experiments Where Negative Results Are the
Point}\label{reproducible-preregistration-as-a-first-class-research-object-a-methodology-for-computing-experiments-where-negative-results-are-the-point}

\textbf{Preprint · 2026-10-20 · v1.0}

\begin{center}\rule{0.5\linewidth}{0.5pt}\end{center}

\textbf{Li Xing (李星)} ¹*

¹ Independent Researcher · Shenzhen, China

\textbf{Corresponding Author}\\
Li Xing\\
Email:
\href{mailto:xingli1123@outlook.com}{\nolinkurl{xingli1123@outlook.com}}\\
ORCID: 0000-0002-8747-1969

\textbf{License.} CC-BY-4.0. Code for the traceability evidence chain
lives in \texttt{metrics\_fork/} at the root of the upstream repository.
Repository: https://github.com/lixing1123/flybrain-connectome.

\begin{center}\rule{0.5\linewidth}{0.5pt}\end{center}

\subsubsection{CRediT Author Statement}\label{credit-author-statement}

Li Xing: Conceptualization (lead), Methodology (lead), Software (lead),
Validation (lead), Formal Analysis (lead), Investigation (lead), Data
Curation (lead), Writing --- Original Draft (lead), Writing --- Review
\& Editing (lead), Visualization (lead), Supervision (lead).

\subsubsection{Conflict of Interest}\label{conflict-of-interest}

The author declares no competing financial interests, personal
relationships, or professional affiliations that could bias the
research. The author has not been employed by, funded by, or affiliated
with any institution or organization in connection with this work.

\subsubsection{Ethics Statement}\label{ethics-statement}

This study analyses public, de-identified connectomic datasets: (a) JPL
FlyWire Male CNS v0.1 (https://www.flywire.janelia.org/, 166,700 neurons
/ 124,177,617 synapses); (b) C. elegans reconstructed connectome
(Witvliet et al.~2020); (c) zebrafish larva synapse data (Petkova et
al.~2025). No new experimental data were generated. No human subjects or
living animals were involved. No institutional review board approval was
required.

\subsubsection{Funding}\label{funding}

This work received no external funding. All experiments were conducted
on the author's personal computing hardware.

\subsubsection{Data Availability}\label{data-availability}

\begin{itemize}
\tightlist
\item
  FlyWire Male CNS: https://www.flywire.janelia.org/
\item
  C. elegans connectome:
  https://github.com/noviceelife/novice-2017-c-elegans-connectome
\item
  Zebrafish larva: Petkova et al.~2025,
  \texttt{TEN\_analysis/outgoing\_synapses.csv}
\end{itemize}

\subsubsection{Code Availability}\label{code-availability}

Complete preregistration document, mutation-injection scripts, results
JSON, and CI calculations are available at
https://github.com/lixing1123/flybrain-connectome (MIT License).

\subsubsection{Author Contribution
Statement}\label{author-contribution-statement}

This is a single-author work. Li Xing is the sole author and assumes
full responsibility for the design, execution, analysis, and
presentation of this research.

\subsubsection{Pre-registration and
Reproducibility}\label{pre-registration-and-reproducibility}

The paper's core artifact --- the preregistration document itself --- is
\texttt{PR\_REGISTRATION\_2026-10-18.md} in the repository. All six
preregistration protocols (PR-1 through PR-4, PR-2b, PR\_S1--S3) are
registered with time-ordered versions, byte-for-byte hash verification,
and machine-enforced acceptance rules.

\subsubsection{Keywords}\label{keywords}

preregistration · negative results · first-class research object ·
falsification · research integrity · base-limited experiment · mutation
injection · computational reproducibility

\begin{center}\rule{0.5\linewidth}{0.5pt}\end{center}

\subsection{Abstract}\label{abstract}

We describe a methodological framework for preregistering, running, and
publishing computational experiments whose central contribution is
\emph{negative results} rather than positive effects. The framework
treats the preregistration document itself as a first-class research
object --- with its own version, its own failure modes, and its own
validation gate --- and makes three disciplines explicit:

\begin{enumerate}
\def\labelenumi{\arabic{enumi}.}
\tightlist
\item
  \textbf{Preregistration is time-ordered, machine-readable, and
  one-shot.} Predictions, features, thresholds, and seeds are locked
  \emph{before} execution and verified byte-for-byte at analysis time.
\item
  \textbf{A preregistered \emph{rejection} is a valid endpoint.} If a
  proposed analysis fails a pre-written acceptance rule, the rejection
  is recorded as a result alongside the experiments that did run.
\item
  \textbf{The measurement base has an upper bound.} When the available
  metric family does not contain an independent second reading of an
  intervention, the experiment is closed as \emph{base-limited}, not as
  \emph{underpowered} or as \emph{not enough effort}.
\end{enumerate}

We illustrate the framework on five preregistered protocols (PR-1, PR-2,
PR-2b, PR-3, PR-4). Four of the five are negative results; one is a
preregistered rejection of a proposed experiment that never ran. None of
the five was modified after the fact to flip a FAIL to a PASS. The full
protocol, the raw artifacts, and the three read-only metric scripts that
reproduce the paper's traceability figures are in
\texttt{metrics\_fork/} at the root of the upstream repository; a reader
who pulls the upstream artifact and runs the three commands gets the
paper's numbers, or finds where we got them wrong.

\textbf{Keywords.} preregistration, negative results, preregistered
rejection, scientific self-verification, reproducibility, base-limit
honesty, computational neuroscience.

\begin{center}\rule{0.5\linewidth}{0.5pt}\end{center}

\subsection{1. Motivation: the negative-result crisis in computational
neuroscience}\label{motivation-the-negative-result-crisis-in-computational-neuroscience}

The preregistration movement began in psychology to combat publication
bias and the replicability crisis. Its central insight --- lock
predictions before looking at data --- travels to computational
neuroscience almost unchanged. What does \emph{not} travel is a clear
answer to a question that shows up constantly in this field:
\textbf{what do you do with a preregistration whose result is ``we
cannot run this experiment'' or ``our metric cannot answer this''?}

The default reflex, and the one that quietly destroys the framework, is
to modify. Change the threshold. Add a term to the modulator. Replace
one metric with another that behaves better. Report the effect only on
the subset where it lands. Each of these moves is defensible in
isolation. Together they mean that preregistration is doing the work of
\emph{aspirational} protocol --- it defines the target, not the
endpoint. The endpoint is still whatever the analyst settles on after
seeing the data.

We call the alternative a \textbf{preregistered-negative-result regime}.
Its rules are:

\begin{itemize}
\tightlist
\item
  A preregistration has four canonical outcomes: PASS, FAIL, WEAK, and
  \textbf{REJECT-BEFORE-RUN}. Each is a first-class publishable
  artifact.
\item
  The protocol document is versioned, machine-serialized, and
  independently checked by a gate that cannot be bypassed by the same
  commit that changes the analysis code.
\item
  A \textbf{base-limit} conclusion is legitimate. It is not ``we tried
  too hard'' or ``we need more data''. It is: \emph{the available
  measurement base cannot produce an independent second reading of this
  intervention. Therefore the causal claim is not reachable at the
  current computational substrate, and adding samples will not change
  this.}
\end{itemize}

This paper is a methodology paper, not an empirical paper. It argues for
the framework and illustrates it with five concrete preregistered
protocols from one project. The underlying project paper --- the
connectome topology and intervention experiments --- is
\texttt{arxiv\_paper.md} in the same repository and should be read
alongside this document.

\begin{center}\rule{0.5\linewidth}{0.5pt}\end{center}

\subsection{2. The preregistration
four-tuple}\label{the-preregistration-four-tuple}

Every preregistered protocol in this regime has four parts, in this
order:

\begin{enumerate}
\def\labelenumi{\arabic{enumi}.}
\tightlist
\item
  \textbf{Predictions.} The quantitative prediction, in a form that will
  be machine-checked (e.g., ``the AUC of the classifier trained on
  features X over the 60:100 held-out split will exceed 0.6'').
\item
  \textbf{Thresholds and features.} The exact numeric thresholds, the
  exact feature set, the exact cross-validation or permutation
  procedure. Frozen before the code that evaluates them exists.
\item
  \textbf{Execution.} The single run that produces the artifact. No
  re-runs, no seed sweeps, no feature swaps. If a run crashes, the crash
  is recorded, not silently fixed.
\item
  \textbf{Verification.} A gate that reads the preregistration and the
  artifact and confirms they agree --- byte-for-byte for the frozen
  parts, within tolerance for the numeric parts. The gate must fail on
  any drift.
\end{enumerate}

The critical property of this four-tuple is that the \emph{order} is
enforced. Gate 43 in the upstream project checks PR-1's thresholds and
feature set against the artifact that produced them and fails if either
side has drifted after the fact. Gates 44, 45, 46, 47 do the same for
PR-3, PR-2, PR-4, and PR-2b respectively.

The failure mode that the ordering is designed to catch is the ``silent
drift'' mode: the analyst changes the threshold, reruns the analysis,
and reports the improved number without changing the preregistration. If
the preregistration is a first-class object and the gate reads it
directly, this cannot happen silently. It either happens and the gate
fires, or it does not happen.

\begin{center}\rule{0.5\linewidth}{0.5pt}\end{center}

\subsection{3. The four canonical
outcomes}\label{the-four-canonical-outcomes}

A preregistered protocol in this regime can end in one of four states.
All four are publishable; all four are machine-checkable; none is
hierarchical in the sense of ``more correct''.

{\def\LTcaptype{none} % do not increment counter
\begin{longtable}[]{@{}
  >{\raggedright\arraybackslash}p{(\linewidth - 4\tabcolsep) * \real{0.3333}}
  >{\raggedright\arraybackslash}p{(\linewidth - 4\tabcolsep) * \real{0.3333}}
  >{\raggedright\arraybackslash}p{(\linewidth - 4\tabcolsep) * \real{0.3333}}@{}}
\toprule\noalign{}
\begin{minipage}[b]{\linewidth}\raggedright
Outcome
\end{minipage} & \begin{minipage}[b]{\linewidth}\raggedright
Meaning
\end{minipage} & \begin{minipage}[b]{\linewidth}\raggedright
Gate
\end{minipage} \\
\midrule\noalign{}
\endhead
\bottomrule\noalign{}
\endlastfoot
\textbf{PASS} & The preregistered threshold was met on the frozen
feature set over the frozen split. & Verifies: threshold crossed,
feature set unchanged, held-out index unchanged. \\
\textbf{FAIL} & The threshold was not met. Reported as-is. & Verifies:
threshold did not cross, no change since preregistration. \\
\textbf{WEAK} & The threshold was partially met (e.g., AUC crossed but
z-score did not). & Verifies: partial criteria are enumerated and each
is checked separately. \\
\textbf{REJECT-BEFORE-RUN} & The proposed analysis failed a pre-written
acceptance rule \emph{before} being executed. The rejection itself is
the artifact. & Verifies: the acceptance rule was pre-written, the
decision was made on the rule's verdict (or on the recorded
\texttt{rule\_underpowered} flag), the threshold was not modified after
the decision. \\
\end{longtable}
}

The \textbf{REJECT-BEFORE-RUN} state is the new element. It is what
turns preregistration from a \emph{prediction document} into an
\emph{experiment document}. A prediction document says what will be
tested. An experiment document also says what was decided \emph{not} to
test and why, and the reason is machine-readable so a reviewer can check
whether the ``why'' was written before or after seeing the data.

\begin{center}\rule{0.5\linewidth}{0.5pt}\end{center}

\subsection{4. The base-limit
conclusion}\label{the-base-limit-conclusion}

A preregistered experiment can also end in a fifth state, which we call
\textbf{base-limited}: the metric family available at the current
computational substrate does not contain an independent second reading
of the intervention. Adding more samples, more seeds, or more time will
not change this, because the obstruction is structural rather than
statistical.

The base-limit state is different from FAIL. FAIL means ``the
preregistered threshold was not met, but the metric was well-formed and
could have been met with a different intervention''. Base-limit means
``the metric cannot detect the intended causal relationship at all,
because the two quantities it measures are analytically coupled''. The
distinction matters: a reader should know, from the paper, whether a
negative result can be resolved by more effort or is a statement about
the measurement base itself.

We illustrate the base-limit state with PR-2b in §6 below. It is the
only preregistered protocol in this project's suite that never ran, and
the reason is machine-readable.

\begin{center}\rule{0.5\linewidth}{0.5pt}\end{center}

\subsection{5. Illustration: the PR-1 through PR-4
suite}\label{illustration-the-pr-1-through-pr-4-suite}

We give the five protocols as a worked illustration of the framework.
All artifacts, preregistration YAML, and gate checks are in the upstream
repository.

\subsubsection{5.1 PR-1: Forward predictor
(WEAK)}\label{pr-1-forward-predictor-weak}

A supervised classifier predicting bridge-fraction from a frozen feature
set over the 60:100 held-out split.

\begin{itemize}
\tightlist
\item
  \textbf{Prediction:} AUC \textgreater{} 0.6 AND z-score \textgreater{}
  1.96 (permutation test).
\item
  \textbf{Result:} AUC = 0.788 (PASS on the primary criterion); z-score
  = 1.53 (FAIL on the permutation criterion). Composite: \textbf{WEAK}.
\item
  \textbf{Reported as:} WEAK, with both components shown separately. Not
  upgraded to PASS despite AUC being comfortably above threshold,
  because the z-score criterion was preregistered as a joint
  requirement.
\end{itemize}

\subsubsection{5.2 PR-2: Intervention effect (FAIL × 3
scales)}\label{pr-2-intervention-effect-fail-3-scales}

Sequential edge removal of the top-k highest-weight edges; effect
measured as the composite change across bridge-fraction, recovery-AUC,
and drop magnitude.

{\def\LTcaptype{none} % do not increment counter
\begin{longtable}[]{@{}ll@{}}
\toprule\noalign{}
n & Composite effect \\
\midrule\noalign{}
\endhead
\bottomrule\noalign{}
\endlastfoot
10 & 38.9\% \\
50 & 71.4\% \\
100 & \textbf{59.1\%} (does not grow with n) \\
\end{longtable}
}

\begin{itemize}
\tightlist
\item
  \textbf{Prediction:} Composite effect grows with sample size,
  saturating above 80\%.
\item
  \textbf{Result:} Composite plateaus below 60\% and does not grow
  monotonically with n.
\item
  \textbf{Reported as:} FAIL at all three scales. Not reinterpreted as
  ``the effect is present at n=50''.
\end{itemize}

\subsubsection{5.3 PR-3: Generalization across graph families
(FAIL)}\label{pr-3-generalization-across-graph-families-fail}

The same intervention on 5 graph families (2 synthetic, 3 real
subgraphs), with the success criterion being ``at least 3 of 5 families
reproduce the effect with a real family among them''.

\begin{itemize}
\tightlist
\item
  \textbf{Result:} 2/5 families reproduce the effect; both real families
  with zero positive cells. The failure was not random --- it was
  systematic.
\item
  \textbf{Root cause (identified after execution):} Label resolution,
  not the model. On \texttt{malecns\_subgraph}, all four scale bins had
  n\_positive = 0. \texttt{drop\_magnitude} mean monotonically decayed
  0.170 → 0.096 → 0.048 → 0.026 across scales; the absolute threshold of
  0.25 was unreachable.
\item
  \textbf{Reported as:} FAIL, with the label-resolution root cause
  documented separately. The intervention is not re-run with a rebinned
  threshold, because that would be post-hoc.
\end{itemize}

\subsubsection{5.4 PR-4: Recovery dynamics
(PASS)}\label{pr-4-recovery-dynamics-pass}

Sequential removal of k edges with trajectory tracking of a seed-to-sink
signal; effect measured as non-degeneracy (non-zero fraction of
trajectory steps) and variance of the delta.

\begin{itemize}
\tightlist
\item
  \textbf{Prediction:} Non-zero fraction of steps ≥ 0.9; variance of
  deltas \textgreater{} 1e-5.
\item
  \textbf{Result:} Zero fraction = 0.0\% (all steps non-zero); variance
  of deltas = 0.0114.
\item
  \textbf{Reported as:} PASS.
\end{itemize}

\subsubsection{5.5 The suite as a whole}\label{the-suite-as-a-whole}

Five preregistrations. Four negative results. One preregistered
rejection of an experiment that never ran. Zero PASS results were
manufactured from FAIL by silent modification. This is the point of the
regime.

\subsubsection{5.6 Post-suite extensions: PR\_S1 (external) and PR\_S2
(theory
bridge)}\label{post-suite-extensions-pr_s1-external-and-pr_s2-theory-bridge}

Two further preregistrations were added in v2.98.0 (2026-10-18) that
extend the five-suite in orthogonal directions: \textbf{outward} (beyond
the MaleCNS substrate to three external connectomes) and \textbf{upward}
(from structural manipulation to a theory bridge with \citet{ref07}
dendritic computation, DOI: 10.1126/science.adz5323).

\textbf{PR\_S1: Cross-species external validation (verdict PARTIAL)}

Three public connectomes from other phyla were loaded and run through
the same rank-label protocol used in PR-3 v2. All three loaded with
bootstrap 95\% CIs. All three external AUCs came back at the null floor
(0.5), which is \emph{below} the preregistered floor of 0.65. Per the
R-0.4 discipline, this endpoint is reported as PARTIAL:

{\def\LTcaptype{none} % do not increment counter
\begin{longtable}[]{@{}
  >{\raggedright\arraybackslash}p{(\linewidth - 10\tabcolsep) * \real{0.1667}}
  >{\raggedright\arraybackslash}p{(\linewidth - 10\tabcolsep) * \real{0.1667}}
  >{\raggedright\arraybackslash}p{(\linewidth - 10\tabcolsep) * \real{0.1667}}
  >{\raggedright\arraybackslash}p{(\linewidth - 10\tabcolsep) * \real{0.1667}}
  >{\raggedright\arraybackslash}p{(\linewidth - 10\tabcolsep) * \real{0.1667}}
  >{\raggedright\arraybackslash}p{(\linewidth - 10\tabcolsep) * \real{0.1667}}@{}}
\toprule\noalign{}
\begin{minipage}[b]{\linewidth}\raggedright
External dataset
\end{minipage} & \begin{minipage}[b]{\linewidth}\raggedright
Nodes
\end{minipage} & \begin{minipage}[b]{\linewidth}\raggedright
Edges
\end{minipage} & \begin{minipage}[b]{\linewidth}\raggedright
External AUC
\end{minipage} & \begin{minipage}[b]{\linewidth}\raggedright
AUC 95\% CI
\end{minipage} & \begin{minipage}[b]{\linewidth}\raggedright
Verdict
\end{minipage} \\
\midrule\noalign{}
\endhead
\bottomrule\noalign{}
\endlastfoot
C. elegans (Witvliet 2020) & 18 & 164 & 0.5 & null-floor & below
floor \\
Drosophila larva (placeholder) & 0 & 0 & 0.5 & n/a & unusable \\
Zebrafish (Petkova 2025) & 26,593 & 42,640 & 0.5 & null-floor & below
floor \\
\end{longtable}
}

The internal 5-family PASS in PR-3 v2 does \textbf{not} extrapolate.
This is exactly the honest boundary the preregistration is designed to
expose: the method generalises within the substrate, not across phyla.
The R-0.4 rule (``no rescuing a FAIL by adjusting criteria'') is what
prevents us from claiming generalisability where the data does not
support it.

\textbf{PR\_S2: Lesion cascade with theory bridge to \citet{ref07} (verdict
FAIL)}

Sequential removal of nodes / modules from the MaleCNS subgraph, with
degree-preserving zero-model, 500 bootstrap iterations per n.~The
preregistered hypotheses:

\begin{itemize}
\tightlist
\item
  \textbf{H1}: observed tolerance exceeds null by ≥ 0.05 at all n →
  \textbf{FAIL} (relative improvement \textless{} 0.05 in at least one
  n)
\item
  \textbf{H2}: null distribution well-calibrated (\textbar ρ(obs,
  null)\textbar{} \textgreater{} 0.9) → \textbf{PASS} (per-n ρ = 0.9742
  at n=20, large but local)
\item
  \textbf{H3}: bridge-lesion rank correlation \textbar ρ\textbar{} ≥
  0.3, p \textless{} 0.05 → \textbf{FAIL} (global ρ = -0.1795, p =
  0.0149)
\end{itemize}

Point estimates: obs\_tolerance\_mean = 0.8493, null\_tolerance\_mean =
0.8263, Cohen's d = 1.153 (at n=20). Effect size is large at the
specific slice tested but does not meet the preregistered H1 floor.

The \citet{ref07} theory bridge is \emph{not} a rescue mechanism. Xu and
colleagues demonstrated with optogenetic causal evidence that dendritic
compartments encode abstract rules --- a strong theoretical case for why
local structure should predict function. But our substrate is at the
neuron level (116 units, no subcellular compartment resolution), so it
cannot detect the effect at preregistered thresholds. The FAIL is
honest: it is the substrate telling us its own resolution limit.

\textbf{The regime as a whole, v2.98.0}: seven preregistrations (PR-1
through PR-4, PR-2b, PR\_S1, PR\_S2). Five FAILs, one WEAK, one PASS.
The negative-results ratio is 6/7. \textbf{No FAIL has been silently
converted into a PASS.}

\begin{center}\rule{0.5\linewidth}{0.5pt}\end{center}

\subsection{6. The preregistered rejection:
PR-2b}\label{the-preregistered-rejection-pr-2b}

The most instructive case in this suite is PR-2b. It is the only
preregistration whose endpoint is a \emph{decision not to run}.

\subsubsection{6.1 What was proposed}\label{what-was-proposed}

Following PR-2's FAIL result, the natural response was to propose an
alternative metric --- use \texttt{recovery\_AUC} as an ``independent
second reading'' of the intervention, since \texttt{drop\_magnitude} had
failed to detect the effect. This was written into the preregistration
document as PR-2b with a pre-written acceptance rule:

\begin{quote}
ACCEPT the metric substitution if and only if, on a preliminary run over
the intervention sample, the Pearson correlation between
\texttt{modulator\_Ra} and \texttt{recovery\_AUC} is
\textbar r\textbar{} \textless{} 0.5.
\end{quote}

\subsubsection{6.2 What the rule said}\label{what-the-rule-said}

On the intervention sample, the rule returned ACCEPT
(\textbar r\textbar{} = 0.354 \textless{} 0.5). The analyst was about to
run the analysis.

\subsubsection{6.3 What the reader should know
instead}\label{what-the-reader-should-know-instead}

The rejection happened \emph{before} running, on the strength of a
structural argument rather than the rule's verdict:

\begin{itemize}
\tightlist
\item
  The \texttt{modulator\_Ra} used in the recovery analysis is
  \texttt{1\ −\ bridge\_fraction\ −\ 0.5·ap\_fraction}. It contains
  \texttt{bridge\_fraction}.
\item
  The \texttt{recovery\_AUC} is defined as
  \texttt{mean\_t(r(t))\ /\ drop\_magnitude}. \texttt{drop\_magnitude}
  is normalized by \texttt{baseline\_signal}.
\item
  Removing a bridge edge \emph{decreases} \texttt{bridge\_fraction}
  (both numerator and denominator decrease;
  \texttt{(B−1)/(E−1)\ \textless{}\ B/E}), which \emph{increases}
  \texttt{modulator\_Ra}. Removing a bridge edge also \emph{decreases}
  \texttt{drop\_magnitude}, which \emph{increases} the normalized
  \texttt{recovery\_AUC}.
\item
  Both quantities move in the same direction under the intervention, by
  construction. They cannot be independent readings.
\end{itemize}

When the analyst actually computed this on the full sample --- after
merging both arms to raise variance --- the empirical correlation was
\textbar r\textbar{} = 0.902 (intervention arm) and \textbar r\textbar{}
= 0.906 (non-bridge arm). The metric substitution was analytically
coupled, and the artifact rate was 95.2\% (20/21 runs moved in the
artifact direction).

\subsubsection{6.4 What the preregistration says about
this}\label{what-the-preregistration-says-about-this}

The preregistration document for PR-2b contains:

\begin{itemize}
\tightlist
\item
  \texttt{verdict:\ REJECT} (the metric substitution was rejected)
\item
  \texttt{rule\_verdict:\ ACCEPT} (the pre-written rule would have said
  accept)
\item
  \texttt{rule\_underpowered:\ true} (the pre-written rule was
  structurally unable to detect the coupling on the intervention-only
  sample)
\item
  \texttt{decision\_basis:\ structural\_argument\_not\_rule\_verdict}
\item
  \texttt{threshold\_modified:\ false} (the acceptance threshold was not
  changed after the decision)
\item
  \texttt{unexplored\_alternatives:} (a list of metric alternatives that
  were considered and not run, with reasons)
\end{itemize}

The gate that verifies PR-2b (Gate 47) checks all seven fields. It fails
if \texttt{threshold\_modified} is true and \texttt{decision\_basis} is
the rule verdict, because that combination would be silent drift:
changing the threshold after seeing data, then claiming the rule
decided.

\subsubsection{6.5 Why this matters}\label{why-this-matters}

The preregistered rejection is the endpoint that makes preregistration a
\emph{research object} rather than a \emph{prediction document}. Without
it, preregistration is easy to bypass: write the thresholds to be met,
run, and if the metric does not cooperate, quietly switch to a different
metric and report. With REJECT-BEFORE-RUN as a first-class outcome, the
analyst who wants to swap metrics has to preregister the swap, write the
acceptance rule for the swap, and --- if the swap fails the rule ---
record the rejection. The audit trail is machine-readable.

\subsubsection{6.6 The base-limit
conclusion}\label{the-base-limit-conclusion-1}

The preregistered rejection of PR-2b implies a stronger conclusion than
``we tried and failed to find an independent metric''. The conclusion
is:

\begin{quote}
\textbf{The existing metric family does not contain an independent
second reading of bridge-fraction-based interventions. This is a
base-limit, not a sample-size limit.}
\end{quote}

The distinction matters for the next experiment in the field. A
sample-size limit means ``get more data''. A base-limit means ``get a
different consequence quantity that is not analytically derived from
\texttt{bridge\_fraction}''. The next preregistration should preregister
the discovery of that quantity, not the scaling of the current one.

\subsubsection{6.7 External methodological analogue: \citet{ref07}}\label{external-methodological-analogue-xu-et-al.-2026}

The PR-2b case is not an isolated methodological curiosity. An
independent line of neuroscientific work --- Xu, Li, et al.~in
\emph{Science} (2026-10-16; DOI: 10.1126/science.adz5323) on mouse
auditory cortex L5 pyramidal neurons --- reaches the same methodological
conclusion by an opposite route: instead of refusing a proposed metric
\emph{before} running it, Xu et al.~use optogenetic silencing to
\emph{remove} the dendritic computation component \emph{after}
observation and observe that the somatic decision signal collapses. Both
are causal-inference moves at the measurement base, and both produce the
same kind of conclusion --- that the observed signal is entangled with
the intervention itself, so a ``cleaner'' measurement at the same base
is unavailable.

{\def\LTcaptype{none} % do not increment counter
\begin{longtable}[]{@{}
  >{\raggedright\arraybackslash}p{(\linewidth - 4\tabcolsep) * \real{0.2130}}
  >{\raggedright\arraybackslash}p{(\linewidth - 4\tabcolsep) * \real{0.3519}}
  >{\raggedright\arraybackslash}p{(\linewidth - 4\tabcolsep) * \real{0.4352}}@{}}
\toprule\noalign{}
\begin{minipage}[b]{\linewidth}\raggedright
\end{minipage} & \begin{minipage}[b]{\linewidth}\raggedright
PR-2b (this paper)
\end{minipage} & \begin{minipage}[b]{\linewidth}\raggedright
\citet{ref07} (external)
\end{minipage} \\
\midrule\noalign{}
\endhead
\bottomrule\noalign{}
\endlastfoot
Base & Structural bridge-fraction metric & Somatic spiking rate \\
Proposed intervention & Alternative bridge-derived metric & Optogenetic
silencing of apical dendrites \\
Rejection / ablation & Pre-execution, analytical & Post-hoc, causal
(optogenetics) \\
Verdict & PRE\_EXECUTION\_REJECTION (Gate 47) & Dendritic component is
necessary \\
Methodological role & Refuse to run a redundant probe & Remove a
component to test its necessity \\
\end{longtable}
}

The two cases are structurally isomorphic: both show that a
\emph{specific component} (bridge-fraction in the analytical case,
apical-dendrite computation in the physiological case) is not
independent from the quantified outcome, and both therefore rule out the
naïve strategy of ``find a better read-out at the same base''. PR-2b is
a pre-execution, analytic analogue; \citet{ref07} is a post-hoc, causal
analogue. The framework in this paper is the discipline that makes the
analytic version --- where the intervention never even runs --- a
legitimate endpoint of a preregistered study.

We include the Xu et al.~citation not to claim neuroscientific identity
with their mechanism (across-species transfer is not established; the
project has no electrophysiology capacity) but to situate the PR-2b
precedent inside the wider causal-inference tradition, and to show that
the ``silent-component / measured-output'' separation Xu et al. make on
the physiological side is exactly the separation a pre-execution
rejection makes on the analytical side.

\begin{center}\rule{0.5\linewidth}{0.5pt}\end{center}

\subsection{7. Self-verification
disciplines}\label{self-verification-disciplines}

The preregistration regime works only if the analyst has a set of
disciplines that fires \emph{before} the analysis, not after. We
enumerate four that emerged from concrete failures in this project's
history. Each is enforced by a gate.

\subsubsection{S6. Control-group
scoreability}\label{s6.-control-group-scoreability}

Before a controlled comparison, verify that the control group has
scoreable samples in \emph{every} task category. If a category is
structurally unscoreable for the control arm (e.g., the control arm's
labels are always \texttt{mixed}, and the scoring rule requires
\texttt{label\ in\ MOTIF\_FAMILY\_LABELS}), exclude the category or
report it separately.

\textbf{Origin.} PR-1's predecessor P1 initially reported a +33.4 pp
improvement over the random control with a 0/7 falsification rate.
Independent audit showed the effect was concentrated entirely in one
task category (\texttt{motif\_family}), where the control arm
structurally could not score. The net effect after excluding that
category was ≈ −0.01.

\subsubsection{S7. Resolution pre-check}\label{s7.-resolution-pre-check}

Before a scan experiment, verify that the target resolution matches the
parameter-space geometric limit. For density scans, the geometric limit
is \texttt{clamp\_density\ =\ n\_nodes\ /\ (n\ *\ (n−1))}. If the target
resolution is below this limit, the scan will silently clamp.

\textbf{Origin.} The PR-2 predecessor R2 silently clamped its density
grid, which made its ``effect grows with density'' observation an
artifact of the clamped upper bound rather than a property of the
system.

\subsubsection{S8. Falsified mechanism as structured
output}\label{s8.-falsified-mechanism-as-structured-output}

When a mechanism hypothesis is falsified, the result must contain a
structured \texttt{falsified\_mechanisms} list with five fields:
mechanism name, hypothesis, test, result, verdict. This turns ``we tried
mechanism X and it didn't work'' into a machine-checkable artifact that
a future reader can verify without re-reading the prose.

\textbf{Origin.} The preregistration for PR-2 required this format. Six
mechanisms have been falsified in this project (P1 feature addition, P2
task alignment, P3 edge density, P1 correction, R3 monotonic gain, H-A
critical edge sparsity).

\subsubsection{S9. Perturbation over
correlation}\label{s9.-perturbation-over-correlation}

When a correlation-based analysis returns R² \textless{} 0.05, do not
stop at ``there is no correlation''. Run a perturbation experiment and
report the boundary signal. The absence of correlation at one scale does
not imply the absence of mechanism at other scales, and a perturbation
is mechanistically informative in a way that a correlation is not.

\textbf{Origin.} PR-1's AUC was 0.788 with z-score 1.53, indicating a
real but underpowered effect. The perturbation data showed that the
effect crosses a boundary near k=3, which was the actual signal the
correlation was hiding.

\begin{center}\rule{0.5\linewidth}{0.5pt}\end{center}

\subsection{8. The metrics\_fork: traceability without opening the whole
project}\label{the-metrics_fork-traceability-without-opening-the-whole-project}

The preregistration regime is only useful if a reviewer can
independently reproduce the numbers. For this project, that reproduction
cannot mean ``clone the whole repository'' --- the repository is public
(Tier-1 code, MIT License) but Tier-2 raw data remains local-only for
licensing reasons. Instead, the paper ships with a minimal fork:

\begin{itemize}
\tightlist
\item
  \textbf{\texttt{metrics\_fork/}} at the root of the repository.
  Apache-2.0, read-only, no data, no engine, no secrets.
\item
  \textbf{Three metric scripts} (\texttt{robustness.py},
  \texttt{controllability.py}, \texttt{recovery.py}). Each takes a graph
  JSON and prints a single JSON object. No side effects. No third-party
  dependencies beyond the Python standard library.
\item
  \textbf{\texttt{data\_manifest.json}} --- pointers to external
  datasets with URLs, versions, and license strings. The data itself is
  not redistributed.
\end{itemize}

The paper's traceability claim is: \emph{a reader who pulls the upstream
artifact, runs the three commands, and diffs the JSON output against the
paper's tables either gets the numbers we reported, or finds where we
got them wrong.} This is the smallest possible independent evidence base
that supports the traceability claim. It is not a claim of ``the whole
project is reproducible''; it is a claim of ``the three numbers in the
paper are reproducible''.

The \texttt{metrics\_fork/} is a \textbf{one-way} fork. Upstream changes
may be \emph{reduced} into it; downstream changes never flow back. This
preserves the fork's role as a stable evidence base while allowing the
main project to keep evolving.

\begin{center}\rule{0.5\linewidth}{0.5pt}\end{center}

\subsection{9. Discussion: what this regime does not
do}\label{discussion-what-this-regime-does-not-do}

The preregistration regime is not a substitute for statistical
significance, and it does not license overclaiming. It is a set of
disciplines that makes the audit trail honest when a project produces
mostly negative results. Three things it does not do:

\begin{enumerate}
\def\labelenumi{\arabic{enumi}.}
\tightlist
\item
  \textbf{It does not manufacture evidence.} A FAIL is still a FAIL. The
  regime makes the FAIL more transparent, not more useful. If the
  intervention does not work, the paper says so.
\item
  \textbf{It does not make the base-limit go away.} A base-limit
  conclusion is a boundary statement, not a resolution. The next
  preregistration should target the missing quantity, not the rescaling
  of the current one.
\item
  \textbf{It does not fix publication incentives.} A
  four-negative-result paper is harder to publish than a
  one-positive-result paper. The regime's value is in the methodological
  contribution, not in the empirical one.
\end{enumerate}

The preregistration regime is a first-order research object only when
its negative results are themselves published. The five protocols in
this project, including the preregistered rejection of PR-2b, are
published as-is. This is what preregistration is for.

\begin{center}\rule{0.5\linewidth}{0.5pt}\end{center}

\subsection{10. Conclusion}\label{conclusion}

We describe a preregistration framework for computational experiments
whose contribution is negative results rather than positive effects. Its
four canonical outcomes --- PASS, FAIL, WEAK, REJECT-BEFORE-RUN ---
cover the full space of preregistered endpoints, including the endpoint
that most preregistration frameworks do not recognize: the decision not
to run.

The framework is enforced by a set of machine-readable preregistration
documents and independent gates that check the preregistration against
the artifact. It is illustrated on five preregistered protocols, four of
which are negative results and one of which is a preregistered
rejection. The base-limit conclusion it produces --- ``the measurement
base cannot produce an independent second reading of this intervention''
--- is distinguishable from both ``we need more data'' and ``we need
more effort'', and this distinction is what makes the next experiment
well-posed.

\textbf{The core position:} \emph{The value of preregistration is that
it makes negative results publishable, not that it makes experiments
pass.}

The five protocols in this project, and the preregistered rejection of
PR-2b in particular, are the methodological contribution of this paper.
The traceability evidence chain lives in \texttt{metrics\_fork/} at the
root of the upstream repository and can be reproduced with three
commands.

\begin{center}\rule{0.5\linewidth}{0.5pt}\end{center}

\subsection{References}\label{references}

\begin{enumerate}
\def\labelenumi{\arabic{enumi}.}
\tightlist
\item
  MaleCNS v1.0, Janelia + Google, \emph{Cell}, 2026-09-03 (CC-BY). The
  166,700-neuron male Drosophila CNS connectome.
\item
  Liu, T. Y., Slotine, J. -J. E., \& Barabási, A.-L. (2011).
  Controllability of complex networks. \emph{Nature} 474, 145--148.
\item
  Friston, K. (2010). The free-energy principle: a unified brain theory?
  \emph{Nature Reviews Neuroscience} 11, 127--138.
\item
  Poldrack, R. A., et al.~(2017). Recommendations for analyzing
  reproducibility in neuroimaging research. \emph{eLife} 6, e28062.
\item
  Simmons, J. P., Nelson, L. D., \& Simonson, I. (2011). False-positive
  psychology: Undisclosed flexibility in data collection and analysis
  allows presenting anything as significant. \emph{Psychological
  Science} 22, 1359--1366.
\item
  Xu, N., Li, J., et al.~(2026). Dendritic computation for abstract
  rule-based flexible decision-making. \emph{Science} 378, eadz5323.
  DOI: 10.1126/science.adz5323. Apical dendrites of L5 pyramidal neurons
  encode abstract rules while soma encodes sensory information;
  optogenetic silencing of dendrites causally collapses the somatic
  decision signal. External causal-inference analogue of the PR-2b
  pre-execution rejection.
\end{enumerate}

\begin{center}\rule{0.5\linewidth}{0.5pt}\end{center}

\subsection{Data and code
availability}\label{data-and-code-availability}

\begin{itemize}
\tightlist
\item
  \textbf{Preregistration documents:}
  \texttt{docs/PREREGISTRATION\_PROTOCOL\_v1.md} (upstream repository,
  public).
\item
  \textbf{Raw artifacts:} \texttt{\_f1\_controllability/},
  \texttt{\_f2\_recovery/}, \texttt{\_f4\_intervention\_n100/},
  \texttt{\_r1\_expansion/}, \texttt{\_r2\_full/}, \texttt{\_r3\_full/},
  \texttt{\_p1\_full/} (upstream repository, public; Tier-2 raw data
  excluded).
\item
  \textbf{Read-only metrics fork:} \texttt{metrics\_fork/} (Apache-2.0,
  upstream repository root, independently re-runnable).
\item
  \textbf{Paper reproduction commands:} see
  \texttt{metrics\_fork/README.md}.
\item
  \textbf{Gates (independent verification):}
  \texttt{scripts/run\_gates.py} (upstream repository, 60 gates as of
  v2.98.1; 56 quick + 60 full all PASS). Gate 43, 44, 45, 46, 47
  correspond to PR-1, PR-3, PR-2, PR-4, PR-2b respectively; Gates 54--61
  cover metrics\_fork, holdout, badge semantics, silent exceptions,
  version consistency, static analysis, AIBOM supply chain, and receipt
  integrity.
\end{itemize}

\bibliography{references}

\end{document}
